\section{What has to be settled before claiming}\label{sec:open}

\subsection{The checks this review could not close}

\begin{enumerate}[label=Q\arabic*.,leftmargin=*]
\item \textbf{Term-by-term comparison with \citet{Yang2026}.} Contribution C2
      asserts that the closure-specific closed forms are markedly simpler than
      the general forest-simplex algebra, and C3 asserts that $O(n)$ pricing is
      available here without the structural overhead of Yang's model. Neither
      has been verified against the paper's actual formulas. This is the single
      most load-bearing open item: if Yang's development specialises to
      something equivalent, C2 and C3 both collapse into ``an instance of a
      known scheme''.
\item \textbf{Exact delimitation against \citet{MollakhaniEtAl2026}.} Their
      abstract settles the direction: the algorithm is motivated by the
      structure of basic feasible solutions of the simplex method, and it then
      merges and splits chunks directly on the dependency graph, maintaining a
      partition ordered by decreasing aggregate ratio rather than an LP basis.
      That is a departure from the simplex, not a specialisation of it. What
      remains unchecked is the full text: whether the paper anywhere derives
      reduced costs, a ratio test or dual multipliers for \eqref{eq:class}, and
      how far its optimality argument coincides with the pricing argument used
      here.
\item \textbf{The network simplex with one side row.} \emph{Closed on the
      structural question; open only on the term-by-term comparison.}
      \citet{ChenSaigal1977} and \citet{Caliskan2011} have now been obtained
      and read. They push this block back to 1977 and settle what it does and
      does not cover.

      \citet{ChenSaigal1977} treat a capacitated flow problem with $k$
      arbitrary additional linear rows. Their Proposition 1.1.5 is the
      structural statement: in any partition of a basis, the square block
      corresponds naturally to a \emph{spanning tree} of the network, so a basis
      is a spanning tree plus $k$ further columns, and the subnetwork carried by
      the basis has rank $n-1$ and nullity $k$. Pivoting, pricing and inversion
      are then driven by a \emph{working basis} of size $k$ rather than $n+k-1$.
      \citet{Caliskan2011} is the case $k=1$ made explicit on maximum flow with
      a budget row: a basis is the two subtrees rooted at $s$ and $t$, the arc
      $(t,s)$, and \emph{one} further arc closing a fundamental cycle, with his
      Lemma~1 recording that the net cost around that cycle cannot vanish, which
      is the rank condition the side row imposes.

      Two consequences. First, R7 in \cref{sec:retracted} is confirmed and can
      now be cited precisely instead of being called textbook: tree-plus-$k$
      bases, a working basis of the size of the side rows, and a specialised
      pricing built on them are due to 1977, not to 2017. Second, the phrase
      ``to the best of our knowledge'' in \cref{sec:positioning} survives, for a
      reason \cref{sec:ratio} already states: this whole block works on the
      node--arc incidence matrix with variables on \emph{arcs}, where a basis is
      a spanning tree, while \eqref{eq:class} carries variables on
      \emph{vertices} and difference constraints on arcs --- the
      \emph{transposed} incidence matrix --- where a basis is a forest with
      anchored roots and not a spanning tree of the same graph. The paradigm is
      theirs; the specialisation to the transposed structure is not in any of
      the three papers.

      What remains open is narrower and is the same check as Q1:
      \citet{HolzhauserEtAl2017}, the closest of the block, has been obtained
      but not read term by term, so the claim that the closure-specific closed
      forms are simpler than the generic working-basis algebra is still asserted
      rather than demonstrated for that paper.
\item \textbf{Formal separation from the normalized tree.} This is now as
      load-bearing as Q1. A normalized tree of \citet{LerchsGrossmann1965} and
      an \FS{} basis are, on inspection, the same object read twice: forest
      plus root, aggregates per branch, a sign test per branch, one arc in and
      one root arc out. What has to be written down is the exact statement of
      the correspondence and of where it breaks --- presumably that a normalized
      tree is not a basis of the ratio LP, that it carries no ratio test and no
      dual multiplier, and that the parametric Lerchs--Grossmann of
      \citet{Hochbaum2001} walks breakpoints rather than performing simplex
      pivots on the normalised program. If instead the correspondence turns out
      to be exact, then \FS{} is a reinterpretation of a 1965 algorithm, which
      is a publishable observation but a completely different paper from the one
      currently drafted.
\item \textbf{Prior art on the equality normalisation.} \emph{Closed, and it
      is a precedent.} \citet{CharnesCooper1962} has now been obtained and
      read. Their Theorem~2 reduces a linear fractional program to the ordinary
      linear program
      \[
        \max\ c^{\mathsf T}y + \alpha t
        \quad\text{s.t.}\quad
        Ay - bt \le 0,\qquad
        d^{\mathsf T}y + \beta t = 1,\qquad y,t \ge 0,
      \]
      in which the normalisation is an \emph{equality} and the constant is
      exactly $1$. The program \eqref{eq:class} that \FS{} solves is the
      homogeneous case $\alpha = \beta = 0$, $d = w$, of that transformation.
      The deliberate choice of an equality is therefore not a modelling
      contribution of this project; it is the 1962 construction, and
      \cref{sec:ratio} already cited it as such.

      Their treatment of sign is worth recording because it is \emph{not} ours.
      Charnes and Cooper require $\operatorname{sgn}(\gamma) =
      \operatorname{sgn}(d^{\mathsf T}x^\ast + \beta)$ and therefore solve two
      linear programs, the second obtained by negating $(c,\alpha)$ and
      $(d,\beta)$ --- ``a change in sign in the functional and in one
      constraint''. Here the weights are positive, so the denominator is
      positive on the whole feasible cone and only their first program is ever
      needed, while a negative \emph{numerator} is representable inside it.
      That is an observation about the instance class, not a device, and it
      should not be presented as one.
\end{enumerate}

\subsection{The computational axis}

The project's own experimental evidence does not support a performance claim.
On the frozen test bed the max-flow backends --- parametric decomposition and
Dinkelbach --- solve more instances, and solve them faster, than \FS{} in every
configuration tried; the numbers are in the experimental report and are not
reproduced here, because they are being re-measured on a quiet machine and this
document must not carry a second, drifting copy of them. The qualitative
finding is stable across every campaign run so far and is unlikely to move: the
dominant cost is degeneracy, the primal spends the overwhelming majority of its
pivots on zero steps, and the dual removes the degeneracy without buying a
comparable speedup, trading null steps for small ones.

The consequence for positioning is that the work is, today, a
\emph{specialisation} paper and not a \emph{performance} paper, and the three
experiments that could change that have not been run:

\begin{enumerate}[label=E\arabic*.,leftmargin=*]
\item \textbf{Single ratio oracle.} \FS{} against a spanning-forest
      linearization implementation adapted to the first-chunk task, against a
      static min-cut/root search, and against a general LP solver.
\item \textbf{Full decomposition path.} Repeated residual oracles against
      spanning-forest linearization, against parametric max-flow
      \citep{GalloEtAl1989}, against the parametric Lerchs--Grossmann of
      \citet{Hochbaum2001}, and against a mining-style implementation
      \citep{ByunDimitrakopoulos2013} where available.
\item \textbf{Sequences of related instances.} Vertex deletions, forced
      inclusions and exclusions, residual prefix removal, small profit and
      weight perturbations, arc additions and removals --- comparing forest
      basis reuse against pseudoflow warm start, against parametric state
      reuse, and against cold restart.
\end{enumerate}

E2 was the natural place to look for an advantage, and it has now been measured
on the full canonical sequence rather than on a single capacity. The result is
negative and the reason is instructive, so both are recorded here.

The task \texttt{full-path} was run on twelve instances spanning the generated
families, a historical PCKP instance and a MineLib instance, comparing the tuned
primal, the dual, Dinkelbach and the parametric decomposition, two replicas
each, under a sixty-second limit. The parametric backend is fastest on eleven of
the twelve; the dual wins only on the arc-free graph, where there is nothing to
pivot on. The margin is not marginal: the median ratio between the tuned primal
and the parametric backend is a factor of twelve, and on a layered DAG with $n =
2000$ it is a factor of $859$ --- sixty seconds against seven hundredths. On the
largest instance tried, a sparse DAG with $n = 20{,}000$, the tuned primal, the
dual and Dinkelbach all exhaust the limit while the parametric backend finishes
the whole sequence in $3.6$ seconds.

The diagnosis matters more than the ranking. The oracle count is exactly what
the design predicts --- one ratio oracle per closure layer, the same count
Dinkelbach pays --- and the parametric backend uses about two flow computations
per layer. The difference is entirely inside the oracle: each one costs tens to
hundreds of thousands of pivots, of which between $98\%$ and $100\%$ are
degenerate. On the layered instance the primal spends $768{,}639$ pivots, essentially
all of them zero steps, to produce a sequence the parametric backend obtains
with a few hundred flow computations.

The warm start was then implemented and measured, and it recovers most of that
cost. The extracted block \emph{is} the positive component, that is one tree of
the basis forest, so the other trees survive untouched and stay anchored: the
next oracle can reuse that forest instead of rebuilding from the empty basis.
The condition that was missing from the project's earlier decision is a single
one and is checked rather than proved: the values are $1/w(T)$ on the positive
tree $T$ and zero elsewhere, so an active arc leaving $T$ would violate
feasibility and $T$ must be closed in the residual. The best closed surviving
tree is chosen, and if none is closed the oracle restarts cold. The check costs
$\mathcal{O}(m)$ per oracle.

The effect is on the pivot count, which is where the time was: $251{,}780$
pivots become $3{,}080$ on a disconnected DAG, $89{,}468$ become $6{,}108$ on a
sparse one, $48{,}974$ become $6{,}915$ on \texttt{minelib\_newman1}. The
decomposition produced is identical block by block and the objective of a
capacity query is unchanged to nine decimals. In wall time the warm primal is
$3.1$ times faster than the cold one, overtakes Dinkelbach, and wins three of
the ten instances every configuration solves --- the families whose canonical
sequence is longest --- against one of twelve before.

This changes the size of the gap, not its direction. The parametric backend
still wins the other seven instances and is $3.4$ times faster in total, so the
positioning of \cref{sec:positioning} stands: this is a specialisation paper,
not a performance paper. What has changed is that the axis is no longer
hypothetical. The degeneracy is reduced, not removed, and two measurements are
still owed: the frozen configuration was tuned on single-capacity solves and
never on the full decomposition --- a paired ablation on this task puts a sink
initial basis $1.21$ times ahead of it --- and the anti-cycling regime has never
been tuned for this task at all.

An end-to-end
branch-and-cut comparison comes after that, not before, and must not be used as
the first piece of evidence.

\subsection{Honest summary}

For the forest structure, the basic-solution characterisation, the merge/split
pivots, the sign test on branch aggregates, maximum-ratio closure, the canonical
decomposition and the incidence-matrix-plus-side-row paradigm, \FS{} is not new
--- and the forest itself is not new since 1965, not since 2025. What may be new is
the specialisation of the ordinary simplex, primal and dual, to this particular
linear program, the closed forms that specialisation produces, and the use of it
to drive the entire canonical sequence of blocks rather than the first ---
and that is a narrow claim which \cref{sec:positioning} states narrowly, and
which Q1, Q2 and Q4 have to close before it is made in public --- Q3 and Q5
having been closed by the readings recorded above, both of them against the
claim rather than for it.
